\movedhead{The continuous measure}
\label{sec:pctmeasure}\label{app:pctmeasure}
Amendment 3 added the continuous measure because the categorical scale was at
ceiling in the pilot, and it earned its place. It finds 22 non-zero cells
against the categorical measure's 10, and the gain is concentrated exactly
where the categorical scale has least room. The four-point scale is
effectively three-point in practice, because \texttt{BUY\_MORE} was never
once returned by any model at any effort level, and each model concentrates
on a single category: \texttt{SELL\_PARTIAL} on 80.8\% of Opus responses,
\texttt{HOLD} on 63.2\% of Sonnet's, and a near-even split between those two
in Haiku. Opus is the extreme case. Its categorical answer barely moves
between conditions, so the categorical measure finds 3, 0 and 0 non-zero
cells at \texttt{low}, \texttt{medium} and \texttt{high}, while the graded
percentage inside that one category moves freely and finds 2, 6 and 5. Two of
the four Holm surviving cells exist only on the continuous measure. This is
the clearest vindication of Amendment 3: on the registered measure alone,
the strongest model would have appeared almost bias-free at every effort
level, and that appearance would have been an artifact of a scale with no
room in it.

The two measures agree on sign and ordering wherever both move, and they
disagree on individual responses rarely and systematically. Eighty-six of
9{,}720 responses, 0.88\%, pair a categorical answer with a contradictory
percentage. Every one of the 86 is Opus; 84 of them are at \texttt{low}
effort, a rate of 2.59\% against 0.06\% at \texttt{high}; and every one is a
\texttt{HOLD} accompanied by a non-zero percentage to sell. These rows are
flagged and kept on both measures, never reconciled, as registered. They also
account for the one cell that breaks the disposition sign pattern: Opus at
\texttt{low} effort is the only cell of nine with a negative disposition
effect, and it is the cell in which the model's two answers are least
coherent with each other.
